Modern applications running on the Internet are increasingly distributed to enable scalability, improve service availability, and decrease user latency. In a distributed system, individual components can only coordinate via message passing over a network, making it naturally volatile and prone to partial failures due to subsystems being largely independent \cite{Coulouris-DistributedSystems-2005}. This failure model and the pessimistically asynchronous nature of communication networks, preventing a global notion of time, make many conceptually simple problems notoriously difficult. One example is commonly referred to as the \textit{consensus} problem \cite{Coulouris-DistributedSystems-2005, Fischer-DistributedConsensus-1985, Dwork-ConsensusPartialSynchrony-1988}, and requires multiple processes to decide on the same value in a strongly consistent manner.

Research on this problem has come a long way since the \textit{Paxos} algorithm for solving consensus \cite{Lamport-Paxos-2001} was designed. \textit{Fast Paxos}, for instance, introduced an optimization saving a one-way delay in case there are no concurrent proposals \cite{Lamport-FastPaxos-2006, Junqueira-ClassicVsFastPaxos-2007, Zhao-FastPaxos-2015}. In practise, these algorithms are commonly extended to maintain a continuously growing \textit{distributed log} instead of just being used to decide on a single value. Such a distributed log enables the implementation of \textit{replicated state machines} for processing deterministic operations on multiple nodes in the same manner, providing guaranteed convergence to the same state \cite{Schneider-ReplicatedStateMachines-1990}. This is an important abstraction for applications like distributed databases \cite{Corbett-GoogleSpanner-2013, Taft-CockroachDB-2020, etcd-Docs} and coordination services \cite{Burrows-Chubby-2006, Hunt-Zookeeper-2010}.

While consensus-based distributed logs inherently provide strong safety guarantees, their poor scalability and low throughput can dominate end-to-end system performance, in particular in geo-replicated deployments \cite{Ryabinin-SwiftPaxos-2024, Stathakopoulou-RSMScalability-2022}. As a result, a large chunk of current research focuses on reducing the overhead of log replication by using the optimistic fast paths originally devised for Fast Paxos. Fast-path algorithms, however, rely on favourable execution conditions, which rarely hold in production systems subject to churn, faults, skewed workloads or even targeted attacks. In such collision-heavy scenarios, the expensive recovery mechanisms necessary to satisfy safety guarantees degrade performance even below that of conservative protocols \cite{Junqueira-ClassicVsFastPaxos-2007}.

To mitigate the cost of fast path mistakes, operators currently have the options to either choose protocols that modify the safety constraints for the distributed log \cite{Park-Curp-2019, Ryabinin-SwiftPaxos-2024} or to gate fast paths based on static heuristics \cite{Dobre-HybridPaxos-2010}. However, changing the safety core is not always an option depending on the scenario, whereas static gating rules leave performance on the table in calm periods and cause instability in turbulent ones.

%Todo: Figure out some more about protocols with different safety constraints or hybrid stuff
%Todo: Explain a bit more about conformal risk control

Therefore, this thesis aims to explore the usage of \textit{conformal risk control} \cite{Shafer-ConformalPrediction-2008, Angelopoulus-ConformalPrediction-2023, Angelopoulus-ConformalRiskControl-2025} to regulate when the system can take the fast path and when it must revert to the conservative path. This approach provides distribution-free, finite sample guarantees on a chosen rate of fast path mistakes, leaving safety invariants unchanged.

\generalExpl{Give a general introduction to the area. (Remember to use appropriate references in this and all other sections.)}

\engExpl{Use the glossaries package to help yourself and your readers.
Add the acronyms and abbreviations to lib/acronyms.tex. Some examples are shown below:
In this thesis, we will examine the use of \glspl{LAN}. In this thesis, we will
assume that \glspl{LAN} include \glspl{WLAN}, such as \gls{WiFi}.}


%\section{Background}
%\label{sec:background}
%\sweExpl{svensk: Bakgrund}

%\generalExpl{Present the background for the area. Set the context for your project – so that your reader can understand both your project and this thesis. (Give detailed background information in Chapter 2 - together with related work.)
%Sometimes it is useful to insert a system diagram here so that the reader
%knows what are the different elements and their relationship to each
%other. This also introduces the names/terms/… that you are going to use
%throughout your thesis (be consistent). This figure will also help you later
%delimit what you are going to do and what others have done or will do.}

%As one can find in RFC 1235\,\cite{ioannidis_coherent_1991} multicast is useful for xxxx. A number of different \glspl{OS} have been used in this work, such as the following \glspl{OS}: UNIX, Linux, Windows, etc. The main focus will be on one \gls{OS}, namely Linux.

\section{Problem}
\label{sec:problem}
\sweExpl{svensk: Problemdefinition eller Frågeställning\\
Lyft fram det ursprungliga problemet om det finns något och definiera därefter
den ingenjörsmässiga erfarenheten eller/och vetenskapen som kan komma ur
projektet. }

Longer problem statement\\
If possible, end this section with a question as a problem statement.

% Research Question
\subsection{Original problem and definition}
\label{sec:researchQuestion}
\sweExpl{Ursprungligt problem och definition}
Some text

\subsection{Scientific and engineering issues}\sweExpl{Vetenskaplig och ingenjörsmässig frågeställning}
some text

\section{Purpose}
\sweExpl{Syfte}
\sweExpl{Skilj på syfte och mål! Syfte är att förändra något till det bättre. I examensarbetet finns ofta två aspekter på detta. Dels vill problemägaren (företaget) få sitt problem löst till det bättre men akademin och ingenjörssamfundet vill också få nya erfarenheter och vetskap. Beskriv ett syfte som tillfredställer båda dessa aspekter.\\
Det finns även ett syfte till som kan vara värt att beakta och det är att du som student skall ta examen och att du måste bevisa, i ditt examensarbete, att du uppfyller examensmålen. Dessa mål sammanfaller med kursmålen för examensarbetskursen. 
}
\generalExpl{State the purpose of your thesis and the purpose of your degree project.\\
Describe who benefits and how they benefit if you achieve your goals. Include anticipated ethical, sustainability, social issues, etc. related to your project. (Return to these in your reflections in Section~\ref{sec:reflections}.)}



\section{Goals}
\sweExpl{Mål}
\sweExpl{Skilj på syfte och mål. Syftet är att åstakomma en förändring i något. Målen är vad som konkret skall göras för att om möjligt uppnå den önskade förändringen (syfte). }

\generalExpl{State the goal/goals of this degree project.}

The goal of this project is XXX. This has been divided into the following three sub-goals:
\begin{enumerate}
\item Subgoal 1 \sweExpl{för att tillfredsställa problemägaren – industrin?}
\item Subgoal 2\sweExpl{för att tillfredsställa ingenjörssamfundet och vetenskapen – akademin) }
\item Subgoal 3\sweExpl{eventuellt, för att uppfylla kursmålen – du som student}
\end{enumerate}

\generalExpl{In addition to presenting the goal(s), you might also state what the deliverables and results of the project are.}

\section{Research Methodology}\sweExpl{Undersökningsmetod}
\sweExpl{Här anger du vilken vilken övergripande undersökningsstrategi eller metod du skall använda för att försöka besvara den akademiska frågeställning och samtidigt lösa det e v ursprungliga problemet. Ofta kan man använda ”lösandet av ursprungsproblemet” som en fallstudie kring en akademisk frågeställning. Du undersöker någon intressant fråga i ”skarpt” läge och samlar resultat och erfarenhet ur detta.\\
Tänk på att företaget ibland måste stå tillbaka i sin önskan och förväntan på projektets resultat till förmån för ny eller kompletterande ingenjörserfarenhet och vetenskap (ditt examensarbete). Det är du som student som bestämmer och löser fördelningen mellan dessa två intressen men se till att alla är informerade. }
\generalExpl{Introduce your choice of methodology/methodologies and method/methods – and the reason why you chose them. Contrast them with and explain why you did not choose other methodologies or methods. (The details of the actual methodology and method you have chosen will be given in Chapter~\ref{ch:methods}. Note that in Chapter~\ref{ch:methods}, the focus could be research strategies, data collection, data analysis, and quality assurance.)\\
In this section you should present your philosophical assumption(s), research method(s), and research approach(es).}

\section{Delimitations}\sweExpl{Avgränsningar}
\generalExpl{Describe the boundary/limits of your thesis project and what you are explicitly not going to do. This will help you bound your efforts – as you have clearly defined what is out of the scope of this thesis project. Explain the delimitations. These are all the things that could affect the study if they were examined and included in the degree project.}

\section{Structure of the thesis}\sweExpl{ Rapportens disposition}
Chapter~\ref{ch:background} presents relevant background information about xxx.  Chapter~\ref{ch:methods} presents the methodology and method used to solve the problem. …